\section{Improved Elastic Weight Consolidation}

The preceding derivation leaves five implementation choices: global penalty strength, damping, aggregation weights, output geometry, and matrix approximation. For old-task samples $a_1,\dots,a_N$, IEWC uses
\[
M_{\mathrm{IEWC}}
=
\frac1N
\sum_{n=1}^N
\bar h_n
\frac{g_ng_n^\top}{r_n^\top G_n^{-1}r_n+\tau_A},
\qquad
g_n=J_n^\top r_n,
\]
where $\bar h_n\ge0$ has empirical mean one. The continual-learning penalty is
\[
\frac{\lambda}{2}
(\theta-\theta_A^\star)^\top
M_{\mathrm{IEWC}}
(\theta-\theta_A^\star).
\]

The canonical geometry is $G_n=I$. We use task-relative damping
\[
\tau_A
=
\kappa\,\operatorname{median}
\{r_i^\top G_i^{-1}r_i:r_i\ne0\},
\]
with $\kappa=10^{-2}$ in the confirmatory comparison. This makes $\kappa$ dimensionless and preserves global loss-unit invariance. If all denominators are zero, the stored importance is zero.

Uniform IEWC sets $\bar h_n=1$. To study weighting without confounding it with normalization, we define a separate \emph{residual-diversity weighting} (RDW) variant. Let
\[
u_n=\frac{g_n}{\|g_n\|_2}
\]
for nonzero parameter gradients. In a fixed sample order, RDW maintains an orthonormal basis $Q$ of at most $K$ earlier residual directions and assigns
\[
h_n=\|(I-Q^\top Q)u_n\|_2^2.
\]
When the residual is nonzero and the cap is not reached, its normalized direction is appended to $Q$; the final $h_n$ are mean-normalized. This streaming rule is motivated by residual-gradient diversity ideas in GSS \citep{GSS}, but it is an order-dependent aggregation heuristic, not part of the level-set derivation. Crucially, both OPG-RDW and IEF-RDW receive the same unit parameter-gradient directions. Scalar IEF normalization cannot change $u_n$, so the $2\times2$ comparison
\[
\{\mathrm{OPG},\mathrm{IEF}\}
\times
\{\mathrm{uniform},\mathrm{RDW}\}
\]
isolates matrix normalization from sample weighting.

For diagonal training penalties, we take the diagonal after aggregation and rescale every nonzero diagonal to have mean one over trainable parameters. This removes only a global factor absorbable into $\lambda$ and lets all regularized methods use the same within-domain $\lambda$ grid. It does not restore the cross-parameter information discarded by diagonalization.

An IEWC instantiation is specified by the five-tuple
\[
(\lambda,\kappa,\Gamma,h,\Pi),
\]
where $\lambda$ sets preservation strength, $\kappa$ sets relative damping, $\Gamma$ maps an old output to its local metric $G_n$, $h$ supplies aggregation weights, and $\Pi$ is the matrix approximation operator (for example, full, diagonal, or low rank). These are genuine design choices rather than values fixed by the derivation. Scalar loss preservation gives OPG-EWC; loss-level-set normal preservation gives uniform IEWC. We call this CLOC-derived mechanism \emph{Improved Elastic Weight Consolidation} while treating its individual axes separately in theory and experiment.
